\subsection{Por qué CV convencional ya no es suficiente}

Una validación convencional toma una partición aleatoria o agrupada de las 138 etiquetas, pero
no intenta que el pseudo-test se parezca específicamente a las 59 parcelas de competencia.
En un problema de objetivo fijo, la pregunta relevante es condicional a $X_{\Target}$.

Checkpoint~04B introduce una pseudo-competición. Para una máscara
$\Target^\star\subset\Train$:
\begin{enumerate}
\item se oculta $y_i$ para $i\in\Target^\star$;
\item se mantiene visible $X_{\Target^\star}$;
\item toda transformación $X$-only puede usar las 197 covariables;
\item cualquier fit supervisado usa sólo
$\Train^\star=\Train\setminus\Target^\star$;
\item se producen predicciones para $\Target^\star$;
\item sólo después se revela $y_{\Target^\star}$ para calcular métricas.
\end{enumerate}

De esta forma la frontera de información de cada pseudo-split reproduce la frontera real:
\[
(X_{\Train^\star},y_{\Train^\star},X_{\Target^\star})
\quad\longrightarrow\quad
\widehat y_{\Target^\star}.
\]

\subsection{Tamaño de la pseudo-competición}

La competencia tiene una fracción objetivo
\[
\pi_U=\frac{59}{197}\approx0.2995.
\]
Aplicada al universo de 138 etiquetas:
\[
138\pi_U\approx41.3.
\]
Por eso cada pseudo-competición target-matched y state-random oculta 41 parcelas.

Se congelaron:
\begin{itemize}
\item 16 target-matched repeats;
\item 8 state-random repeats;
\item 5 folds state-stratified legados;
\item 5 folds municipality-grouped legados.
\end{itemize}
En total se evalúan 34 splits con propósitos distintos.

\subsection{Construcción target-matched}

Las máscaras target-matched no usan $y$. Entre múltiples candidatos se busca aproximar la
geometría del conjunto real en:
\begin{itemize}
\item composición state/municipality;
\item nearest geographic distance;
\item nearest agronomic distance;
\item temporal similarity;
\item adversarial target probability.
\end{itemize}

La calidad se resume, entre otros, por total variation de la distribución municipal y distancia
entre perfiles estandarizados. El procedimiento logra:
\[
TV_{\mathrm{municipio}}:
0.0547\ \text{target-matched}
\quad\text{vs}\quad
0.0889\ \text{state-random},
\]
y
\[
d_{\mathrm{perfil}}:
0.5369\ \text{target-matched}
\quad\text{vs}\quad
0.9227\ \text{state-random}.
\]
Esto verifica que el protocolo principal se parece más a los 59 objetivos en $X$, sin
conocer sus $y$.

\safefigure{../../reports/checkpoint_04b/figures/split_match_quality.png}
{Calidad de matching de las pseudo-competiciones respecto al conjunto objetivo real.}
{fig:cp04b-match}

\subsection{Modelos candidatos}

04B mezcla baselines globales y métodos explícitamente locales.

\paragraph{Geographic kNN.}
Para una consulta $q$ y sus $k$ vecinos visibles $\mathcal N_k(q)$:
\[
\hat y_q
=
\frac{\sum_{i\in\mathcal N_k(q)}
(d_{qi}+\epsilon)^{-p}y_i}
{\sum_{i\in\mathcal N_k(q)}
(d_{qi}+\epsilon)^{-p}}.
\]

\paragraph{Local Ridge.}
Para cada consulta se construye un training set diferente de $k$ vecinos y se resuelve
\begin{equation}
\hat\beta_q
=
\arg\min_{\beta}
\sum_{i\in\mathcal N_k(q)}
w_{qi}
(y_i-x_i^\top\beta)^2
+\alpha\|\beta\|_2^2,
\qquad
w_{qi}=(d_{qi}+\epsilon)^{-1}.
\label{eq:local-ridge}
\end{equation}
La predicción es $\hat y_q=x_q^\top\hat\beta_q$. Este modelo es query-specific: dos objetivos
pueden usar subconjuntos de etiquetas diferentes.

\paragraph{Graph Laplacian.}
A partir de una matriz de distancias $D_X$ se crea un kNN simétrico con pesos RBF:
\[
w_{ij}
=
\exp\left[
-\frac12\left(\frac{d_{ij}}{\sigma}\right)^2
\right],
\]
donde $\sigma$ es la mediana de distancias de aristas positivas. Sea $A=(w_{ij})$,
$D=\diag(A\bm 1)$ y $L=D-A$. Con máscara $M$ de nodos etiquetados se resuelve
\begin{equation}
(M+\lambda L+\epsilon I)\bm f
=
M(\bm y-\bar y\bm 1).
\label{eq:graph-laplacian}
\end{equation}
La predicción en nodos query es $\bar y+f_q$. Esta formulación propaga una función suave sobre
la geometría $X$ \citep{Zhou2004,Belkin2006}.

\paragraph{Distancias mixtas.}
Se evaluaron combinaciones normalizadas:
\[
d_{\mathrm{mix}}
=
a\,d_{\mathrm{geo}}
+b\,d_{\mathrm{agro}}
+c\,d_{\mathrm{temporal}},
\qquad a+b+c=1.
\]
El mejor resultado inicial no requiere una componente temporal dominante, consistente con 04A.

\subsection{Resultados target-matched}

\begin{table}[H]
\centering
\caption{Ranking principal de Checkpoint 04B.}
\label{tab:cp04b-primary}
\begin{tabular}{lrrrr}
\toprule
Método & RMSE medio & mediana & peor split & comentario \\
\midrule
LocalRidge $k=20,\alpha=100$ & \textbf{0.4946} & 0.4778 & 0.6927 & mejor media \\
Graph $k=8,\lambda=2$ & 0.4974 & 0.4878 & 0.6995 & robusto \\
LocalRidge $k=30,\alpha=100$ & 0.4976 & 0.4847 & 0.6950 & cercano \\
GeoKNN $k=10$ & 0.5021 & 0.4794 & 0.7137 & baseline local \\
Blend PLS8+Graph8 & 0.5093 & 0.5012 & 0.6968 & blend \\
GlobalPLS C0 4 & 0.5379 & -- & -- & global \\
\bottomrule
\end{tabular}
\end{table}

El mejor local mejora aproximadamente 8\% respecto a GlobalPLS C0 en el mismo protocolo
target-matched:
\[
\frac{0.5379-0.4946}{0.5379}\approx0.0805.
\]
Esto es una ganancia material comparada con las diferencias de milésimas entre variantes
globales de 03C.

\safefigure{../../reports/checkpoint_04b/figures/primary_method_ranking.png}
{Ranking de métodos bajo las 16 pseudo-competiciones target-matched.}
{fig:cp04b-ranking}

\subsection{Robustez del grafo}

El Graph $k=8,\lambda=2$ obtiene:
\[
\begin{array}{lr}
\text{target-matched} & 0.4974\\
\text{state-random} & 0.4655\\
\text{state-stratified} & 0.4779\\
\text{municipality-grouped} & 0.5466.
\end{array}
\]
Su peor familia, 0.5466, es muy inferior al grouped 0.7481 de E123. No es una comparación
perfectamente idéntica porque 04B permite topología transductiva $X$-only; precisamente ése es
el punto: la estructura del problema permite un estimador más apropiado.

LocalRidge lidera en target-matched pero se deteriora mucho más al retirar municipios enteros.
La tensión entre Local y Graph se conserva hasta 04F: Local optimiza la simulación más cercana
a los 59 objetivos, Graph funciona como hedge de robustez espacial.

\subsection{Routing por soporte: primera hipótesis y limitación}

04B agrupó pseudo-targets en tiers low/mid/high según un score estrictamente $X$-only
recalculado dentro del split. En la tabla agregada:
\[
\begin{array}{lll}
\text{low}: & \text{LocalRidge30}, & \RMSE=0.4174,\\
\text{mid}: & \text{GeoKNN10}, & \RMSE=0.5645,\\
\text{high}: & \text{LocalRidge30}, & \RMSE=0.3198.
\end{array}
\]
Esto generó una propuesta inicial para las 59 parcelas: 32 GeoKNN10 y 27 LocalRidge30.

Sin embargo, la misma evidencia pseudo-test se utilizaba para identificar el ganador del tier
y medirlo. Por tanto, el routing es exploratorio. La solución correcta es seleccionar reglas
en splits distintos de aquellos donde se evalúan; 04D.1 introduce precisamente un
leave-one-split-out.

\subsection{Resultado de diseño}

04B cambia el centro del proyecto. Ya no es razonable invertir la mayor parte del esfuerzo en
modelos globales. La evidencia indica:
\begin{enumerate}
\item la geografía transfiere rendimiento;
\item LocalRidge puede superar el mejor anchor global en la geometría target-matched;
\item Graph ofrece una robustez extraordinariamente mejor al extrapolar espacialmente;
\item cualquier routing o tuning adicional debe reutilizar exactamente estas máscaras
congeladas para evitar mover simultáneamente objetivo y método.
\end{enumerate}

El siguiente gran bloque, 04D.1, refina sólo las familias locales/grafo que ya sobrevivieron
esta prueba.
